\setchapterpreamble[u]{\margintoc}
\chapter{Determinants}
\labch{determinants}

\sources{Strang, \emph{Introduction to Linear Algebra}, 6th ed.\ \cite{strang2023ila},
§5.1--5.3; MIT 18.06 \cite{ocw1806}, Lectures~18--20.}

The determinant compresses a square matrix into one number that answers a
yes-or-no question: is the matrix invertible? Its main use in this book is as
the machinery behind eigenvalues in \refch{eigenvalues}.

\section{Three properties that define it}
\labsec{det-properties}

Rather than start from a formula, take three requirements and derive everything
from them. For \(A\in\R^{n\times n}\), the determinant \(\det A\) is the unique
function satisfying:

\begin{enumerate}
	\item \(\det I=1\);
	\item exchanging two rows reverses the sign;
	\item it is linear in each row separately, the other rows being held fixed.
\end{enumerate}

\marginnote{Property~3 is linearity \emph{one row at a time}, not linearity in
the matrix. Multiplying the whole of an \(n\times n\) matrix by \(c\) multiplies
the determinant by \(c^{n}\), once per row.}

\section{What follows immediately}
\labsec{det-consequences}

\begin{proposition}
If two rows of \(A\) are equal then \(\det A=0\).
\end{proposition}

\begin{proof}
Exchanging the two equal rows leaves \(A\) unchanged but reverses the sign by
property~2, so \(\det A=-\det A\).
\end{proof}

\begin{proposition}
Subtracting a multiple of one row from another leaves the determinant unchanged.
\end{proposition}

\begin{proof}
By linearity in the affected row, the determinant of the result is
\(\det A\) minus \(c\) times the determinant of a matrix with two equal rows,
and that second term is zero.
\end{proof}

These two facts make elimination a legitimate way to compute a determinant: the
subtraction steps cost nothing and only the row exchanges have to be counted.

\begin{proposition}
A triangular matrix has determinant equal to the product of its diagonal
entries; a matrix with a zero row has determinant \(0\).
\end{proposition}

\section{Determinant from the pivots}
\labsec{det-pivots}

Combining the last two sections: eliminate \(A\) into \(U\), keeping track of how
many row exchanges were used. Then
\[
	\det A=(-1)^{\,\text{number of exchanges}}\prod_{i=1}^{n}u_{ii},
\]
the product of the pivots up to sign.

\begin{example}
The elimination carried out in \refsec{elimination-example} on
\[
	A=\begin{bmatrix}1&2&1\\3&8&1\\0&4&1\end{bmatrix}
\]
needed no row exchanges and produced the pivots \(1,2,5\), so \(\det A=10\).
\end{example}

\begin{theorem}
\(A\) is invertible if and only if \(\det A\neq0\).
\end{theorem}

\begin{proof}
Elimination produces \(n\) non-zero pivots exactly when \(A\) is non-singular,
and the displayed formula makes \(\det A\) non-zero precisely in that case.
\end{proof}

The multiplicative rule \(\det(AB)=(\det A)(\det B)\) follows from the same
circle of ideas, and gives \(\det(A^{-1})=1/\det A\) as a special case. Transposition also leaves it alone,
\(\det(A\tp)=\det A\), so each statement about rows has a counterpart about
columns.

\section{Cofactors and the big formula}
\labsec{cofactors}

Expanding property~3 across all rows at once gives a sum over the \(n!\)
permutations, one term for each way of picking a single entry from every row and
column. For \(n=3\) this is the familiar six-term expression; for \(n=10\) it has
over three million terms and is useless for computation.

The practical repackaging is cofactor expansion along a row,
\[
	\det A=\sum_{j=1}^{n}a_{ij}C_{ij},
	\qquad
	C_{ij}=(-1)^{i+j}\det M_{ij},
\]
where \(M_{ij}\) deletes row \(i\) and column \(j\).

\begin{example}
Expanding the same matrix along its first row,
\[
	\det A=1(8-4)-2(3-0)+1(12-0)=4-6+12=10,
\]
agreeing with the pivot calculation.
\end{example}

\marginnote{Two routes, one answer. Cofactors are convenient for small or sparse
matrices and for symbolic work; pivots are what any implementation actually
uses.}

\section{Three uses}
\labsec{det-uses}

\paragraph{The inverse formula.}
\(A^{-1}=\dfrac{1}{\det A}\,C\tp\), where \(C\) is the matrix of cofactors. It
makes the dependence of \(A^{-1}\) on the entries of \(A\) explicit, and shows
why a determinant near zero makes the inverse enormous.

\paragraph{Cramer's rule.}
The \(j\)th component of the solution of \(A\vect{x}=\vect{b}\) is
\(x_j=\det B_j/\det A\), where \(B_j\) is \(A\) with column \(j\) replaced by
\(\vect{b}\). Elegant, and far too slow to use.

\paragraph{Volume.}
\(|\det A|\) is the volume of the parallelepiped spanned by the rows of \(A\).
A zero determinant means the solid is flat --- the rows lie in a lower
dimensional subspace --- which is the geometric form of singularity.

\begin{remark}
The volume reading is the one that carries into calculus. In
\refch{jacobian-hessian} the determinant of the Jacobian is the local
volume-scaling factor of a nonlinear map, and it is what appears when a change of
variables is made inside an integral.
\end{remark}
